\section{Round-twenty-seven reconstruction: weak-Harris contraction, typed renewal, and stable orthogonal dynamics}
\label{sec:r27-a4}

The exact remote past is never erased by finitely many append--shift steps.
Accordingly no common finite-step Doeblin minorization is asserted on the
complete-past state.  Mixing is proved in a distance that discounts remote
coordinates, while the operator theory acts on a genuine weighted Lipschitz
space.  The renewal maps are typed before they are composed, and the
Mori--Zwanzig expansion is derived on a declared sectorial coupling domain.

\subsection{Two-rate history geometry}

Let $\mathsf H$, $K(h,dm)$, $U(h,m)$, and the Lyapunov function $W$ be the
objects constructed in A3.  Write
\[
 d_\sigma(h,h')=\sum_{j\ge1}\sigma^j
   (1\wedge d_E(m_{-j},m'_{-j})),\qquad0<\sigma<1.
\]
Stable distortion at depth $j$ is bounded by
$C\vartheta^j\chi(m_{-j})$.  Choose
\[
 0<\vartheta<\rho<1,\qquad
 V(h)=1+\sum_{j\ge1}\rho^j\chi(m_{-j}),\qquad
 W(h)=e^{\epsilon V(h)}.                                  \tag{A4.1}
\]
Then
\[
 PW\le C W^\alpha,\qquad0<\alpha<1,                       \tag{A4.2}
\]
and, on $\{V\le L\}$, changing the past beyond depth $n$ changes the next-mark
kernel by at most
\[
 C_L(\vartheta/\rho)^n                                   \tag{A4.3}
\]
in weighted total variation.  Notice that the exact remote tail itself is
not small and remains a state coordinate.

Define
\[
 d_*(h,h')=\min\left\{1,\,
 \sqrt{d_\sigma(h,h')\,[2+W(h)+W(h')]}\right\}.            \tag{A4.4}
\]
For $0<\gamma<1$ put
\[
 \|f\|_{\mathcal B_\gamma}
 =\sup_h{|f(h)|\over1+W(h)^\gamma}
 +\sup_{h\ne h'}{|f(h)-f(h')|\over d_*(h,h')}.             \tag{A4.5}
\]
The denominator tends to zero with $d_\sigma$, so (A4.5) imposes actual local
continuity.

\begin{lemma}[Local coupling without exact-tail coalescence]
\label{lem:r27-a4-local}
There are $n_0$, $\delta>0$, and $\kappa<1$ such that
\[
 \mathcal W_{d_*}\bigl(P^{n_0}(h,\cdot),P^{n_0}(h',\cdot)\bigr)
 \le\kappa d_*(h,h')                                     \tag{A4.6}
\]
whenever $d_*(h,h')<\delta$.  On every Lyapunov sublevel there is a
$d_*$-small-set estimate
\[
 \sup_{V(h),V(h')\le L}
 \mathcal W_{d_*}\bigl(P^{n_L}(h,\cdot),P^{n_L}(h',\cdot)\bigr)
 \le1-\epsilon_L.                                        \tag{A4.7}
\]
No measure $\nu_L$ is dominated by all exact-tail transition laws.
\end{lemma}

\begin{proof}
For nearby histories, maximally couple the next mark on common regular
branches.  The mismatch probability is controlled by (A4.3) and the common
recent suffix; on a match, $d_\sigma$ contracts by $\sigma$.  The power drift
controls the weight factor in (A4.4), proving (A4.6).  On a Lyapunov sublevel,
truncate only the \emph{influence} of the remote past using (A4.3), cover the
recent regular cylinders by finitely many connector families, and couple the
new finite suffix with fixed positive probability.  The untouched exact tails
may differ, but after $n_L$ shifts their contribution to $d_\sigma$ is
$\sigma^{n_L}$; this proves (A4.7) without impossible exact coalescence.
\end{proof}

\begin{theorem}[Weighted weak-Harris spectral gap]
\label{thm:r27-a4-harris}
There are $C<\infty$ and $\lambda<1$ such that
\[
 \|P^nf-\pi(f)\|_{\mathcal B_\gamma}
 \le C\lambda^n\|f\|_{\mathcal B_\gamma}.                 \tag{A4.8}
\]
The invariant probability is unique among measures with finite $W$ moment,
and the same estimate holds uniformly on the A3 certified compact family.
\end{theorem}

\begin{proof}
Combine the power drift (A4.2), local contraction (A4.6), and sublevel
$d_*$-smallness (A4.7) in the weak-Harris distance.  Iteration gives
exponential Wasserstein contraction.  Duality with point masses controls the
Lipschitz part of (A4.5); the power drift controls the weighted supremum part.
Completeness follows by taking the two limits separately.  Thus the
contraction is an operator estimate on the declared Banach space, not merely
uniqueness of an invariant measure.
\end{proof}

\subsection{Analytic sources on one fixed space}

A source belongs to
\[
 \mathcal U_\delta=\{g:\ |g(h)|\le c_0+\delta\log W(h),\
             \operatorname{Lip}_{d_*}g<\infty\}.
\]
Choose $\delta>0$ and $0<\gamma<1$ so that
$\alpha(\gamma+\delta)<\gamma$.  Define
$P_gf=P(e^gf)$.

\begin{theorem}[Fixed-space Feynman--Kac theorem]
\label{thm:r27-a4-fk}
On a sufficiently small ball of $\mathcal U_\delta$, $P_g$ is a bounded
analytic family on the single space $\mathcal B_\gamma$.  It has an analytic
simple leading eigenvalue, eigenfunction, eigenmeasure, and Doob transform.
\end{theorem}

\begin{proof}
For $|f|\le C(1+W^\gamma)$,
\[
 |P(e^gf)|\le C e^{c_0}P(1+W^{\gamma+\delta})
 \le C'(1+W^{\alpha(\gamma+\delta)})
 \le C''(1+W^\gamma).
\]
The coupling proof with one multiplier difference gives the Lipschitz bound.
The power series in $g$ converges in operator norm because
$(\log W)^k/k!$ is absorbed by $W^\delta$.  The Riesz contour of the isolated
eigenvalue at zero persists on a small source ball.
\end{proof}

\subsection{Typed first/last-return factorization}

Let $\mathcal X$ be the flow anisotropic space and $\mathcal Y=\mathcal
B_\gamma$ the section/history space.  Stable disintegration at a section
point is a probability kernel $\eta_x(dh)$; no noncanonical map $x\mapsto h_x$
is used.  For $\operatorname{Re}z$ large define
\begin{align*}
 B(z)&:\mathcal X\longrightarrow\mathcal Y,\\
 B(z)f(h)&=\int_0^{\tau(m_0(h))}
 e^{-zs}f(\phi_s e^-(m_0(h)))\,ds,\\
 A(z)&:\mathcal Y\longrightarrow\mathcal X,\\
 A(z)F(x)&=\int_{\mathsf H}\!\int_E
 e^{-z\tau(m)}F(U(h,m))\,K(h,dm)\eta_x(dh),\\
 T(z)&:\mathcal Y\longrightarrow\mathcal Y,\\
 T(z)F(h)&=\int_Ee^{-z\tau(m)}F(U(h,m))\,K(h,dm),\\
 R_0(z)&:\mathcal X\longrightarrow\mathcal X.
\end{align*}
The bound (A3.1) and the A2 branch calculus make these operators bounded and
analytic in a common strip.

\begin{proposition}[Typed renewal identity]
\label{prop:r27-a4-renewal}
On the common resolvent set,
\[
 (z-L_{\rm flow})^{-1}
 =R_0(z)+A(z)(I-T(z))^{-1}B(z).                           \tag{A4.9}
\]
Every product in (A4.9) is composable from right to left.
\end{proposition}

\begin{proof}
Apply $B(z)$ to record the entrance segment and its conditional history law.
Each complete return contributes one $T(z)$ on $\mathcal Y$, and $A(z)$
performs the last-return/exit reconstruction in $\mathcal X$.  Fubini is
justified by the exponential roof moment.  Summing
$A(z)T(z)^kB(z)$ gives (A4.9) for large $\operatorname{Re}z$; analytic
continuation gives the identity throughout the common domain.
\end{proof}

\subsection{A proved orthogonal-dynamics gate}

Let $\mathcal H=L^2(\pi)$ and let $L$ be the closed generator of the prepared
continuous-time semigroup.  Let $P$ be a finite-rank physical projection,
$Q=I-P$.  We do not infer stability of $QLQ$ from stability of $L$.
Instead choose the spectral projection $P_0$ onto constants and the finite
pressure tangent modes.  Then $Q_0=I-P_0$ is invariant and
\[
 \|e^{tL}Q_0\|\le Me^{-\gamma_0t}.                        \tag{A4.10}
\]
Assume $P-P_0$ is finite rank and its graph-norm angle obeys
\[
 \|(P-P_0)LQ_0\|+\|Q_0L(P-P_0)\|<\gamma_0/(4M).           \tag{A4.11}
\]

\begin{lemma}[Stable compression theorem]
\label{lem:r27-a4-compression}
On $QD(L)$, $L_Q=QLQ$ generates a semigroup satisfying
\[
 \|e^{tL_Q}Q\|\le2M e^{-\gamma_0t/2}.                     \tag{A4.12}
\]
\end{lemma}

\begin{proof}
Transport $Q\mathcal H$ to $Q_0\mathcal H$ by the finite-rank graph
isomorphism induced by the angle between $P$ and $P_0$.  In those coordinates
$L_Q$ is $L|_{Q_0\mathcal H}$ plus a bounded perturbation whose norm is below
$\gamma_0/(2M)$ by (A4.11).  Duhamel's formula and Gronwall give (A4.12).
This is a separate compression theorem; it never replaces
$e^{tQLQ}$ by $Qe^{tL}Q$.
\end{proof}

\subsection{Sectorial Schur--Feshbach calculus}

Let $\mathcal R=P\mathcal H$.  Require
\[
 B_Q:=QLP\in\mathcal L(\mathcal R,D(L_Q^2)),\qquad
 C_Q:=PLQ\in\mathcal L(D(L_Q),\mathcal R).                \tag{A4.13}
\]
These are verifiable finite-rank coupling conditions.  They imply that
$C_QL_Q^jB_Q$ is bounded for $j=0,1$.

\begin{theorem}[Domain-safe forced memory equation]
\label{thm:r27-a4-memory}
For $x\in D(L)$, $u(t)=Pe^{tL}x$ satisfies
\[
 u'(t)=PLPu(t)+\int_0^tK(t-s)u(s)\,ds+F_x(t),             \tag{A4.14}
\]
where
\[
 K(t)=C_Qe^{tL_Q}B_Q,\qquad
 F_x(t)=C_Qe^{tL_Q}Qx.                                   \tag{A4.15}
\]
The kernel decays exponentially.  For
$\operatorname{Re}z>-\gamma_0/2$,
\[
 \widehat K(z)=C_Q(z-L_Q)^{-1}B_Q,                        \tag{A4.16}
\]
and on every sector $|\arg z|\le\pi-\epsilon$,
\[
 \widehat K(z)
 =z^{-1}C_QB_Q+z^{-2}C_QL_QB_Q+O(|z|^{-3}).              \tag{A4.17}
\]
Moreover,
\[
 P(z-L)^{-1}P
 =\bigl[z-PLP-C_Q(z-L_Q)^{-1}B_Q\bigr]^{-1}.              \tag{A4.18}
\]
\end{theorem}

\begin{proof}
The block equations on $D(L)$ are
$u'=PLPu+C_Qv$ and $v'=B_Qu+L_Qv$.  Variation of constants and Lemma
\ref{lem:r27-a4-compression} give (A4.14)--(A4.16).  For
$y\in D(L_Q^2)$, the resolvent identity
\[
 (z-L_Q)^{-1}y
 =z^{-1}y+z^{-2}L_Qy+z^{-2}(z-L_Q)^{-1}L_Q^2y
\]
and the sectorial resolvent bound prove (A4.17).  Solving the second row of
$(z-L)(u,v)=(f,0)$ and substituting in the first gives (A4.18).
\end{proof}

The leading $z^{-1}$ term is retained.  If an absolutely integrable inversion
remainder is needed, subtract
$(z+\omega)^{-1}C_QB_Q$ and invert that term explicitly; this is bookkeeping,
not a false cancellation.  The active Lyapunov function is exactly the A3
function, recent singularity proximity is included in $\chi$, and all
renewal, source, and memory maps now live on declared domains.
